% Chapter 1: Complete Sentences

\chapter{Complete Sentences}

\section{Pedagogical Purpose}
This chapter lays the foundational understanding of what a sentence is, which is crucial for all subsequent grammar learning. Without a clear grasp of complete sentences, learners cannot effectively build more complex structures, understand parts of speech in context, or write coherently. It establishes the basic unit of communication.

\begin{learningobjectives}
The learner can:
\begin{itemize}
    \item identify complete sentences and sentence fragments.
    \item explain what makes a group of words a complete sentence.
    \item distinguish between a subject and a predicate in a simple sentence.
    \item use capital letters at the beginning of sentences and full stops at the end.
    \item correct sentence fragments by adding missing parts.
    \item create simple, grammatically correct sentences.
\end{itemize}
\end{learningobjectives}

\section{Prerequisite Check}
\begin{quickcheck}
Read the following groups of words. Put a tick (✓) next to the ones that make sense on their own.
\begin{enumerate}
    \item The cat sat on the mat. [(OK)]
    \item Running very fast. [ ]
    \item My friend Maya. [ ]
    \item She likes to read books. [(OK)]
    \item In the garden. [ ]
\end{enumerate}
\end{quickcheck}

\section{Chapter Opener}
\subsection*{A Classroom Conversation}
\textbf{Ms. Sharma:} Good morning, Class 5! Today, we're going to talk about how we put words together to share our thoughts. Rohan, can you tell us what you did this morning?

\textbf{Rohan:} Woke up. Ate breakfast. Ran to school!

\textbf{Ms. Sharma:} That's a great start, Rohan! Maya, what about you?

\textbf{Maya:} I woke up early. Then I read my favourite book. After that, I walked to school with my brother.

\textbf{Ms. Sharma:} Excellent, Maya! Did you notice a difference between how Rohan and Maya spoke?

\textbf{Leo (the parrot, from his cage):} "Woke up! Ran!"

\section{Language Experience}
Read the following short passage.

The sun shone brightly. Birds sang sweet songs. A little squirrel scampered up a tree. Rohan watched it. He smiled. Maya was reading a book. She loved stories. Ms. Sharma wrote on the board. The class was ready to learn.

\section{Look and Notice}
\begin{enumerate}
    \item Which groups of words tell you a complete idea?
    \item What do you notice at the beginning of each of these groups of words?
    \item What do you see at the end of each of these groups of words?
    \item Can you find any groups of words that don't tell you a complete idea?
\end{enumerate}

\section{Discover / Explicit Teaching}
\subsection*{Guided Discovery: What makes a complete thought?}
Let's look at Rohan's words from the opener: "Woke up. Ate breakfast. Ran to school!"

And Maya's words: "I woke up early. Then I read my favourite book. After that, I walked to school with my brother."

Both Rohan and Maya shared what they did. But Maya's sentences felt more complete, didn't they?

A complete sentence needs two main parts:
\begin{enumerate}
    \item \textbf{Who or What} the sentence is about (the \textbf{Subject}).
    \item \textbf{What the Subject does or is} (the \textbf{Predicate}).
\end{enumerate}

Think of it like this:

\textbf{Subject (Who/What) + Predicate (Does/Is)}

For example:
\begin{itemize}
    \item \textbf{The sun} (Who/What) \textbf{shone brightly} (Does/Is).
    \item \textbf{Birds} (Who/What) \textbf{sang sweet songs} (Does/Is).
    \item \textbf{Rohan} (Who/What) \textbf{watched it} (Does/Is).
\end{itemize}

\section{Core Explanation}
A \gramword{sentence} is a group of words that expresses a complete thought. It always has a subject and a predicate.

\begin{examplebox}
    \begin{itemize}
        \item The dog barked loudly.
        \item Maya is a good student.
        \item They play football every day.
    \end{itemize}
\end{examplebox}

Every sentence starts with a \textbf{capital letter} and ends with a \textbf{full stop} (.), a \textbf{question mark} (?), or an \textbf{exclamation mark} (!). In this chapter, we will focus on sentences that end with a full stop.

A complete sentence has two main parts: a \textbf{Subject} and a \textbf{Predicate}.

\begin{examplebox}
    \textbf{The cat} (Subject - Who/What) \textbf{slept soundly} (Predicate - What the cat did).
\end{examplebox}

\begin{warningbox}
    "Slept soundly." (Not a complete sentence, because we don't know \textit{who} or \textit{what} slept soundly.)
    "The cat." (Not a complete sentence, because we don't know \textit{what} the cat did or was.)
\end{warningbox}

\section{Guided Practice}
\begin{activitybox}{Activity A: Identify the Subject and Predicate}
Draw a line under the \textbf{Subject} and circle the \textbf{Predicate} in each sentence.
\begin{enumerate}
    \item \underline{Rohan} \textcircled{likes to play}.
    \item \underline{Maya} \textcircled{reads many books}.
    \item \underline{Ms. Sharma} \textcircled{smiles often}.
    \item \underline{Leo} \textcircled{repeats words}.
    \item \underline{The students} \textcircled{learn grammar}.
\end{enumerate}
\end{activitybox}

\section{Independent Practice}
\begin{activitybox}{Activity B: Sentence or Fragment?}
Write 'S' for Sentence or 'F' for Fragment next to each group of words.
\begin{enumerate}
    \item The little bird. [F]
    \item Flew high in the sky. [F]
    \item The little bird flew high in the sky. [S]
    \item My friend Maya is kind. [S]
    \item Is kind and helpful. [F]
    \item Rohan plays football. [S]
    \item Plays football every day. [F]
    \item The school bell rang. [S]
    \item Rang loudly. [F]
    \item The children laughed. [S]
\end{enumerate}
\end{activitybox}

\section{Summary}
\begin{summarybox}
    \begin{itemize}
        \item A sentence is a group of words that expresses a complete thought.
        \item Every complete sentence has two main parts: a Subject (who or what) and a Predicate (what the subject does or is).
        \item Sentences begin with a capital letter and end with a full stop (.), question mark (?), or exclamation mark (!).
        \item Sentence fragments are incomplete thoughts and are missing a subject or a predicate.
    \end{itemize}
\end{summarybox}